\section{Discussion}\label{sec:discussion}

We close by revisiting the challenges introduced in
Section~\ref{intro}, linking each to the design choices taken in
the methods and knowledge representation sections, situating the
work in the broader landscape of large-scale semantic resources,
and noting outstanding limitations and a sustainability plan.

\subsection{Design choices revisited}

\paragraph{Unified data model.} The first challenge was to represent
heterogeneous observation types under a single vocabulary. We adopted
the Semanticscience Integrated Ontology (SIO) as the upper-level
framework (Section~\ref{sec:knowledge_representation}) precisely
because it offers a small set of relations
(\texttt{hasInput}/\texttt{hasOutput}, \texttt{hasTarget},
\texttt{hasMember}, \texttt{refersTo}, \texttt{hasValue}) that
generalize across sampling, measurement, and provenance patterns.
Sample collection, XRF analysis, DNA extraction, sequencing, and
taxonomic classification are all instantiations of the same
\texttt{Process--hasInput/hasOutput--Object} idiom; quantitative
results are uniformly represented as \texttt{Quantity refersTo
Quality} pairs with explicit units. This uniformity is what allows
cross-domain SPARQL queries to traverse from a climate observation
through a sampling site to a taxon abundance without per-domain
join logic. Recent work on simplified upper-level ontologies, such as SULO
\cite{sulo2025}, takes this argument further: by strengthening a few
fundamental relations and reducing reliance on specialized
predicates, an upper-level ontology can lower the modelling effort
for new domains while preserving interoperability. Our patterns are
compatible with this direction and could be re-grounded in SULO with
limited refactoring.

\paragraph{Taxonomy harmonization.} Section~\ref{sec:methods_taxonomy}
describes our mapping of OTUs from multiple primer/classifier
combinations onto NCBI Taxonomy. Rather than collapsing the
original assertions, we keep both the raw classifier output and
the harmonized identifier, so downstream consumers can re-use the
mapping or substitute their own.

\paragraph{Schema validation at ingest.} We chose ShEx
(Section~\ref{sec:validation}) for ingest validation because the
ETL is batch-oriented and the validator is run from the same
Groovy/Jena toolchain that produces the RDF. SHACL is an obvious
alternative and is supported in-database by several triple stores;
a future port of the shapes from ShEx to SHACL would enable
in-database validation at load time and is a natural extension of
the current pipeline. We deployed Virtuoso as the triple store
because it is available as free software under the GNU General
Public License, supports OWL-Horst forward-chaining
materialization, and has demonstrated operational maturity in
large-scale linked-data deployments.

\paragraph{Semantic query answering at scale.} As reported in
Section~\ref{sec:validation}, full OWL 2 DL reasoning over the
combined TBox and a 37.4M-triple ABox is not computationally
feasible with HermiT on commodity hardware. We therefore restrict
the deductive closure to OWL~2~EL using ELK \cite{elk}. This is a
deliberate trade-off: ELK handles existential restrictions,
conjunction, and role hierarchies, but does not support
disjunction, negation, or inverse roles, so any axioms that rely
on these constructs do not contribute to the materialized closure.
The competency questions in Section~\ref{sec:validation} are
formulated to remain answerable under EL semantics. Two decades of
work on scaling OWL reasoning -- materialization-based approaches
\cite{owlhorst}, query rewriting for the OWL~2~QL profile
\cite{owl2profiles}, modular reasoning, and consequence-based
calculi as implemented in ELK itself -- provide several alternative
points in the design space, and a more expressive reasoner could
be applied to a sub-module of the TBox in future work.

\paragraph{User-facing interface.} The web portal
(Section~\ref{sec:usage}) addresses the challenge that domain
scientists are typically not fluent in SPARQL. The interactive map
and per-site dashboards cover the most common exploratory tasks;
the SPARQL editor exposes the underlying graph for power users,
pre-populated with the competency questions of
Section~\ref{sec:validation} as click-to-load examples.

\subsection{Limitations and outstanding challenges}

Several challenges remain open. First, although ELK closure is
sufficient for the queries described in this paper, more complex
reasoning -- for example over disjunctions arising from ENVO
biome classifications or from negative axioms about taxon
absence -- would require either a different reasoner or the
extraction of a sub-module amenable to a more expressive
calculus. Second, the resource depends on upstream ontologies
(SIO, ChEBI, ENVO, NCBI Taxonomy) that evolve independently;
keeping the deployed graph aligned with their releases is an
ongoing maintenance task and a known reproducibility risk for any
SPARQL query that bottoms out in a specific external IRI. Third,
while the SPARQL editor lowers the barrier to ad-hoc querying,
authoring queries against (alpha-)numeric IRIs is still
demanding; integrating an autocomplete-aware editor such as
YASGUI is a natural next step.

\subsection{Sustainability and updates}

The dataset is versioned and archived on Zenodo, and the
deployment pipeline is reproducible from the public repository.
We intend to extend the resource as additional field expeditions
are processed: new trips will be added as additional ABox
modules without changes to the upper-level patterns. Upstream
ontology updates will be tracked at scheduled intervals;
re-materialization is scripted and runs as part of the
deployment workflow, so an upstream release can be absorbed by
re-running the pipeline. Where an upstream change would alter
the semantics of an existing query, we will record the affected
competency questions in the changelog so that downstream users
can audit the impact.
